\begin{abstract}
Modern omni speech LLMs absorb broad cross-modal, multi-granularity knowledge during pretraining, yet the gain that training-free reinforcement learning can extract at inference time is \emph{not} a property of how powerful, or how ``agentic,'' the system is: it is governed by the \emph{realized reward spread} of the (action-space, reward) pair the practitioner actually optimizes over. We make this precise. Casting training-free RL as Gibbs tilting of a frozen base policy, $\qs \propto \qo\exp(\Rwd/\beta)$, we develop an \emph{optimization-space adequacy} (OSA) account whose qualitative core---and a new strict-positivity lemma---are machine-checked sorry-free in Lean~4, while the quantitative ceilings are Lean theorems that \emph{consume} external concentration inputs (Hoeffding, Pinsker) which we discharge on paper, not inside Lean. \textbf{OSA-1}: the KL-regularized gain obeys $\gain\le\spread^2/(8\beta)$, a bound on the \emph{reward structure}, not the model class. The same frozen model yields vanishing gain under a flat reward---e.g.\ instruction-conditioning of a frozen embedding read-out---and large gain under a high-spread reward---e.g.\ task-level verifiable selection (our own single-model Operator-B: SLURP intent $+0.330$). \textbf{OSA-2}: under context isolation the optimum over isolated blocks \emph{equals} the monolithic optimum on the product reward $\Rwd=\sum_i\Rwd_i$ (\texttt{qstar\_product}), so isolation \emph{per se} adds nothing; the gain is additive (\texttt{gain\_product}) and each \emph{non-degenerate} block contributes \emph{strictly positive} gain (new theorem \texttt{gain\_pos\_of\_nonconstant}), so the only route to more gain is introducing additional non-degenerate verifiable rewards. Whether a frozen speech model affords such per-block rewards is an \emph{empirical} question (Conjecture~1), not a theorem; \textbf{OSA-3} bounds the rollout-stability tax incurred if it does. We test the conjecture and report a deliberately humbling result: the agentic decomposition is empirically supported for \emph{content/intent} but \emph{fails or is fragile for paralinguistics}. A frozen-embedding speaker probe sits near chance ($\approx0.04$) at all $37$ layers, and a rigorous multi-seed re-run of the emotion-pooling gain (dev-selected layer, paired-delta bootstrap, $5$ seeds) gives mean $\Delta=+0.037$ with the paired CI excluding zero in only $2/5$ seeds---the previously reported single-seed $+0.097$ was a winner's-curse artifact of oracle test-layer selection. We instantiate the theory as a self-evolving omni \emph{speech} agent: a generative-omni policy under a verifiable-reward acceptance gate, paired with a memory that uses the omni embedding for \emph{content} and classic speech encoders (ECAPA/SER) for paralinguistics, because the omni embedding does \emph{not} afford load-bearing paralinguistic retrieval. The agent mechanism is not novel; the contributions are the spread-governs-gain theory, the domain transfer, verifiable speech rewards, and the honest paralinguistic negative that bounds where the program works.
\end{abstract}

\section{Introduction}
\label{sec:intro}

A modern omni or speech-capable multimodal LLM is, by the time it is released, an extraordinary repository of latent capability. From large-scale unsupervised audio, parallel speech-text corpora, and instruction mixtures, models such as Qwen2-Audio, Qwen2.5-Omni, Qwen3-Omni, and SALMONN \citep{qwen2audio,qwen25omni,qwen3omni,salmonn2023} acquire transcription, translation, speaker identity, paralinguistic affect, and dialogue competence simultaneously. The animating question of this program is deliberately conservative: \emph{how far can one go without touching a single weight?} We treat the pretrained model as frozen and ask what reward-guided optimization, applied purely at inference time, can elicit. This is the regime of \emph{training-free} reinforcement learning---best-of-$N$ selection, reward-guided decoding, reranking, task-conditioning, and representation read-out---that changes no parameters and no architecture \citep{beirami2024bon,mudgal2024controlled,zuo2025ttrl}.

\paragraph{The lever is the reward structure, not the model class.}
The premise that organizes this paper is that the gain available to inference-time optimization is a property of \emph{which reward, over which action space} one optimizes---not of whether the system is a single model or a multi-component agent. Some reward structures are nearly \emph{flat}: the degrees of freedom they expose---the prompt, the demonstrations, the sampling temperature applied to a fixed read-out---move the model through a narrow envelope of near-identical rewards. In-context demonstrations frequently fix \emph{format} rather than \emph{accuracy}: their labels can be permuted or corrupted with little effect on the answer distribution \citep{min2022rethinking,alice2026}, and for speech in particular, models tend to ``read'' a transcribable surface rather than ``listen'' to the acoustic substrate, so paralinguistic conditioning barely perturbs the output \citep{voxparadox2026,hsu2023iclspeech}. Our own measurements corroborate this for one specific reward---instruction-conditioning of a frozen omni embedding read-out---where the conditional output distribution is close to flat in the reward dimension. But the very same frozen model, optimized against a \emph{different} reward---task-level verifiable selection over generative candidates---exposes a wide reward spread and a large gain. We name the flat case \emph{flat conditioning}; the central object of the theory is the realized reward spread that separates it from the high-spread case, on \emph{one and the same frozen model}.

\paragraph{Training-free RL as Gibbs tilting.}
To reason about this quantitatively we adopt the now-standard variational view of inference-time alignment. Optimizing the KL-regularized objective
\[
\Fobj{\qz} \;=\; \Exp_{z\sim \qz}[\Rwd(z)] \;-\; \beta\,\KL{\qz}{\qo}
\]
over search distributions $\qz$ on the inference-time action space $\cZ$ has the closed-form Gibbs optimum $\qs(z)\propto \qo(z)\exp(\Rwd(z)/\beta)$, with partition function $\Zpart$ \citep{korbak2022kl,beirami2024bon}. Here $\beta$ multiplies the KL penalty and the tilt is $\exp(\Rwd/\beta)$. Every practical training-free method---best-of-$N$ and its soft and asymptotic refinements \citep{verdun2025softbon,yang2024asymptotics,gui2024bonbon}, self-consistency \citep{wang2023selfconsistency}, MBR decoding \citep{ichihara2025mbr,mbrasr2025}, and controlled decoding \citep{mudgal2024controlled}---is an estimator of this tilt over some realization of $\cZ$. The attainable benefit of \emph{any} such estimator is therefore bounded by how far $\qs$ can depart from $\qo$, which in turn is governed by the spread of $\Rwd$ that the base distribution $\qo$ actually places mass on. Flat conditioning is precisely the statement that this realized spread is small \emph{for that reward}---and the same machinery makes ``high-spread reward, large gain'' equally precise.

\paragraph{Central thesis: optimization-space adequacy.}
Our contribution is to turn this observation into a sharp, verifiable account that we call \emph{optimization-space adequacy} (OSA). \textbf{OSA-1 (spread bound on the (action,reward) pair).} The optimization gain is controlled by the realized reward spread: a Hoeffding argument gives $\gain = \Fobj{\qs}-\Fobj{\qo}\le \tfrac{1}{8\beta}\spread^{2}$, where $\spread$ is the range of $\Rwd$ on the support of $\qo$. This ceiling is a statement about the \emph{(action-space, reward) pair}, not the model class. When the spread is small---the flat-conditioning regime---the Gibbs tilt is a vanishing perturbation and more samples, cleverer reranking, and more elaborate search cannot manufacture headroom the conditional distribution does not contain. When the spread is large, the gain can be large \emph{on the very same frozen model}: our own single-model Operator-B task-level selection lifts SLURP intent accuracy by $+0.330$, where flat instruction-conditioning of the same model's embedding read-out moves the factor accuracies by under $0.01$. This mirrors, from the optimization side, the reward-model over-optimization scaling laws observed empirically \citep{gao2023overopt}. \textbf{OSA-2 (isolation is additive but adds nothing by itself).} Decomposing a task across a \emph{context-isolated} system---each isolated context exposing an independent operator on $\cZ$, whether the embedding/Operator-A view $\cZA$ or the generative/Operator-B view $\cZB$---makes the joint action space a product. On a separable product reward $\Rwd(z_1,z_2)=\Rwd_1(z_1)+\Rwd_2(z_2)$ the isolated optimum \emph{equals} the monolithic optimum: $\qs$ factorizes (\texttt{qstar\_product}) and the gain is exactly additive, $\gain=\sum_i\gain_i$ (\texttt{gain\_product}). We are therefore explicit that \emph{isolation per se adds nothing}---it neither creates nor destroys gain on a fixed reward. The \emph{only} new gain comes from introducing additional \emph{non-degenerate} verifiable rewards, and a new sorry-free lemma makes the contribution of each one strict: if a block's reward is non-constant then its gain is strictly positive (\texttt{gain\_pos\_of\_nonconstant}). Whether a given frozen model actually \emph{affords} such non-degenerate per-block rewards---rather than each block collapsing to its own flat band---is an \emph{empirical} question, which we state as Conjecture~1 and test in this paper, not a theorem we prove. \textbf{OSA-3 (the stability tax).} If such blocks exist, optimizing the product action space by rollout incurs variance and feedback-loop pathologies---context collapse, error propagation in append-only memories, and ``memory that hurts'' \citep{ace2025,expfollowing2025,remembering2026}---so the additive gain is realizable only under algorithm-level credit assignment \citep{jitrl2026,ahmadian2024rloo} and an explicit $\beta$-KL trust region that keeps each context's search distribution close to its own base. Together OSA-1/2/3 suggest a design stance: \emph{do not search harder on one model and do not isolate for its own sake; introduce additional non-degenerate verifiable rewards---then measure, per factor, whether the frozen model affords them.}

\paragraph{A self-evolving omni speech agent.}
We instantiate the theory in the most demanding modality. The system pairs (i) a \emph{generative-omni policy} that acts under retrieved context; (ii) a \emph{memory} that is content-addressable over past episodes; and (iii) a \emph{verifiable-reward acceptance gate} that admits a candidate skill or memory only when an automatically checkable speech reward (WER for ASR, verifiable matches for translation and intent, and learned speech reward models for paralinguistics) certifies improvement \citep{tulu3,rlvrincentive2025,gsrm2026,paras2s2025}. We are candid about the memory design in light of our own evidence below: the frozen omni embedding \citep{omniembed2025,omniembedaudio2026} does \emph{not} afford load-bearing paralinguistic retrieval---it is at chance on speaker identity at every layer---so we reserve the omni embedding for \emph{content} and route paralinguistic memory through classic speech encoders (ECAPA-style speaker embeddings, dedicated SER read-outs). The two operators---memory (Operator-A) and policy (Operator-B)---are the context-isolation substrate OSA-2 describes, and the acceptance gate is the trust region OSA-3 demands. To our knowledge no speech-native, no-gradient, self-evolving agent of this form exists; the omni-policy-plus-verifiable-speech-reward pairing is, at present, an unfilled slot, to our knowledge.

\paragraph{What is and is not novel.}
We are explicit and honest about positioning. The \emph{agent mechanism}---reflective self-improvement, skill libraries, memory-augmented policies---is not novel; the general-agent frontier is crowded with strong, closely related designs \citep{reflexion2023,voyager2023,jitrl2026,mem0_2025,memgpt2023}. Our novelty lies elsewhere: in the \emph{optimization-space theory} that shows the gain is governed by realized reward spread rather than model class and pins down exactly what an isolated decomposition can and cannot add; in the \emph{domain transfer} of training-free RL to frozen omni speech models; in \emph{verifiable speech rewards} as the acceptance criterion; and---as a first-class contribution---in an honest \emph{empirical} demarcation of where the program works (content/intent) and where it does not (paralinguistics). The theory is the load-bearing contribution; the system, and the cautionary study around it, are its honest test.

\paragraph{Contributions.}
We make five.
\textbf{(C1) Theory.} We formalize optimization-space adequacy: OSA-1 (spread-bounded gain on the (action,reward) pair), OSA-2 (additivity under isolation, with the isolated optimum equal to the monolithic optimum and each non-degenerate block strictly positive), and OSA-3 (the stability tax). The qualitative core (the gain identity $\gain$, its non-negativity, the flat-reward no-gain case, the product factorizations \texttt{qstar\_product}/\texttt{gain\_product}, the rollout-deficit decomposition) \emph{and} the new strict-positivity lemma \texttt{gain\_pos\_of\_nonconstant} (reward non-constant $\Rightarrow$ $\gain>0$) are sorry-free in Lean~4. The \emph{quantitative} results---the $\spread^2/(8\beta)$ ceiling and the $O(\sqrt{\log N})$ best-of-$N$ regret---are Lean theorems that \emph{consume} their concentration inputs (Hoeffding, Pinsker) as named hypotheses; those inputs are discharged on paper, not inside Lean. We label this scope plainly rather than overstate what the machine has checked.
\textbf{(C2) Convergence analysis.} We connect OSA-3 to a survey of agentic instability (context collapse, append-only error propagation, memory-induced regression) and give conditions, under $\beta$-KL trust regions and credit assignment, for the agent loop to converge rather than plateau or diverge \citep{ace2025,expfollowing2025,remembering2026}.
\textbf{(C3) An honest empirical demarcation, including a cautionary negative.} We report single-model measurements that locate the operating regime and---crucially---we test Conjecture~1 per factor. For \emph{content/intent} the high-spread regime is real: single-model Operator-B selection yields large, well-separated gains ($+0.330$ on SLURP intent, $+0.335$ on URO spoken-QA). For \emph{paralinguistics} the conjecture fails or is fragile, and we say so plainly. A frozen-embedding speaker probe is near chance ($\approx0.04$, chance $0.011$) at all $37$ layers; and a rigorous multi-seed re-run of the emotion-pooling gain---dev-selected layer (no test-set oracle), paired-delta bootstrap, seeds $\{42,7,123,2024,31337\}$---gives mean $\Delta=+0.037$ (range $[-0.053,+0.097]$), with the paired CI excluding zero in only $2/5$ seeds. The previously reported single-seed $+0.097$ was the best seed under \emph{oracle test-layer} selection: a winner's-curse / selection-leak artifact, not a stable effect. We treat this negative as a contribution: it bounds the flagship's paralinguistic claims and is the reason the agent's paralinguistic memory uses classic speech encoders rather than the omni embedding. The re-run is committed at \texttt{projects/speech-mllm-omni-embedding-rl/\_repro/emotion\_pool\_paired\_v2.json}.
\textbf{(C4) System design.} We specify the self-evolving omni speech agent---generative-omni policy, content-addressable memory (omni embedding for content, classic speech encoders for paralinguistics), verifiable-reward gate---as a concrete, reproducible architecture.
\textbf{(C5) Research plan and benchmark.} We lay out a staged research plan and propose a \emph{cross-session paralinguistic memory} benchmark for speech agents, addressing a gap in current memory and multi-turn audio evaluations \citep{longmemeval2024,locomo2024,audiomc2025}.

\paragraph{Organization.}
\Cref{sec:intro} has framed the spread-governs-gain thesis and the OSA account. The remainder of the paper develops the formal setting and the Gibbs-tilt objective; states OSA-1 (the spread bound on the action--reward pair), OSA-2 (additivity under isolation, the equality of the isolated and monolithic optima, the strict-positivity lemma, and the empirical Conjecture~1 it rests on), and OSA-3 (the stability tax), with an explicit account of which Lean inequalities are sorry-free and which consume named external inputs; analyzes convergence of the agent loop under trust-region credit assignment; presents the self-evolving omni speech agent and its verifiable-reward gate together with feasibility, the single-model results, and the paralinguistic cautionary study; introduces the proposed cross-session paralinguistic memory benchmark; and closes with related work, limitations, and the staged research plan.